% Generated by book/tools/notebook_to_tex.py. Do not edit by hand.

\chapter[벤치마크 평가 (Benchmark Evaluation)]{분야별 벤치마크 평가 (Benchmark Evaluation)}

\label{ch:29}

\markboth{벤치마크 평가 (Benchmark Evaluation)}{벤치마크 평가 (Benchmark Evaluation)}

\chaptermeta{29_benchmark_eval/29_benchmark_eval.ipynb}{https://colab.research.google.com/github/yoon-gu/neuqes-101/blob/master/29_benchmark_eval/29_benchmark_eval.ipynb}{생성형 LLM 평가가 분류 평가와 다른 이유를 task format, log-likelihood, few-shot, LLM-as-judge 관점에서 정리}



\index{1Benchmark Evaluation@Benchmark Evaluation}

\index{1LLM evaluation@LLM evaluation}

\index{1KoBEST@KoBEST}

\index{1Qwen2 5 0 5B Instruct@Qwen2.5-0.5B-Instruct}

\index{1multiple choice@multiple choice}

\index{1log likelihood@log-likelihood}

\index{1acc norm@acc\_norm}

\index{1exact match@exact match}

\index{1generation evaluation@generation evaluation}

\index{1answer extraction@answer extraction}

\index{1few shot@few-shot}

\index{1zero shot@zero-shot}

\index{1in context learning@in-context learning}

\index{1lm evaluation harness@lm-evaluation-harness}

\index{1lm eval@lm-eval}

\index{1LLM as judge@LLM-as-judge}

\index{1human evaluation@human evaluation}

\index{1Goodhart s law@Goodhart's law}

\index{1contamination@contamination}

\index{1leaderboard@leaderboard}

\index{0벤치마크 평가@벤치마크 평가}

\index{0생성형 평가@생성형 평가}

\index{0객관식 평가@객관식 평가}

\index{0로그우도@로그우도}

\index{0길이 정규화@길이 정규화}

\index{0정답 추출@정답 추출}

\index{0퓨샷@퓨샷}

\index{0제로샷@제로샷}

\index{0문맥 내 학습@문맥 내 학습}

\index{0LLM 심판@LLM 심판}

\index{0사람 평가@사람 평가}

\index{0벤치마크 오염@벤치마크 오염}

\index{0리더보드@리더보드}

\index{1KoBEST HellaSwag@KoBEST HellaSwag}

\index{1KoBEST BoolQ@KoBEST BoolQ}

\index{1HFLM@HFLM}

\index{1simple evaluate@simple\_evaluate}

\index{1TaskManager@TaskManager}

\index{1loglikelihood@loglikelihood}

\index{1acc@acc}

\index{1stderr@stderr}

\index{1MMLU@MMLU}

\index{1KMMLU@KMMLU}

\index{1GSM8K@GSM8K}

\index{1HumanEval@HumanEval}

\index{1MT Bench@MT-Bench}

\index{1LogicKor@LogicKor}

\index{1Chatbot Arena@Chatbot Arena}

\index{1position bias@position bias}

\index{1verbosity bias@verbosity bias}

\index{1judge ceiling@judge ceiling}

\index{1gold standard@gold standard}

\index{1EXAONE@EXAONE}

\index{1Gemma@Gemma}

\index{1Qwen@Qwen}

\index{1GLM@GLM}

\index{1DeepSeek@DeepSeek}

\index{1tech report@tech report}

\index{0평가 생태계@평가 생태계}

\index{0평가 도구@평가 도구}

\index{0평가 함정@평가 함정}

\index{0정확도@정확도}

\index{0정규식 추출@정규식 추출}

\index{0판정 편향@판정 편향}

\index{0전문가 평가@전문가 평가}

\index{0대중 평가@대중 평가}

\index{0평가 항해 전략@평가 항해 전략}



\section{누적 추적표}

\begin{booktable}{29장 누적 추적표 표}{tab:ch29-01}
\begin{adjustbox}{max width=\textwidth}
\begin{tabular}{@{}llllll@{}}
\toprule
Ch & 모델 & 단계 & 데이터 & 학습 축 & 평가 \\
\midrule

27 & KoGPT2 \inlinecode{skt/kogpt2-base-v2} (125M) & continual pretraining & 한국어 TinyStories 30K & next-token CE & perplexity / 생성 샘플 \\
28 & KoGPT2 (125M, SFT) & SFT (instruction tuning) & KoAlpaca instruction-response & next-token CE (response-only) & instruction following (정성) \\
\textbf{29 \(\leftarrow\) 여기} & \textbf{Qwen2.5-0.5B-Instruct (+KoGPT2 SFT 대조)} & \textbf{(학습 없음 --- 평가만)} & \textbf{KoBEST / 산술 subset} & \textbf{--- (평가 장)} & \textbf{task format 별: MC log-likelihood / 생성+추출 / LLM-judge} \\
30 (다음) & \ref{ch:28}장 SFT 모델 + ref & alignment (DPO) & preference (chosen/rejected) & DPO loss (response-only) & preference 정렬 \\
\bottomrule
\end{tabular}
\end{adjustbox}
\end{booktable}

\textbf{평가 장이라 학습 축이 없습니다.} \ref{ch:24}--\ref{ch:28}장이 \emph{능력을 만드는} 단계였다면, 본 장은 \emph{그 능력을 측정} 하는 단계입니다. \ref{ch:30}장부터는 다시 학습 (alignment) 으로 돌아갑니다.

전체 장 표는 \href{https://github.com/yoon-gu/neuqes-101\#챕터별-변화추적표}{루트 README} 를 참고하세요.

\begin{center}\rule{0.5\linewidth}{0.5pt}\end{center}

\section{\texorpdfstring{GPT 시대 학습 4단계 --- 본 장은 \emph{측정} 단계}{GPT 시대 학습 4단계 --- 본 장은 측정 단계}}

\begin{booktable}{29장 GPT 시대 학습 4단계 - 본 장은 측정 단계 표}{tab:ch29-02}
\begin{adjustbox}{max width=\textwidth}
\begin{tabular}{@{}llll@{}}
\toprule
단계 & 용어 & 본 커리큘럼 & 본 장? \\
\midrule

1 & Pretraining & \ref{ch:24}장 (영어), \ref{ch:26}장 (한국어) & \\
2 & Continual pretraining & \ref{ch:25}장 (영어), \ref{ch:27}장 (한국어) & \\
3 & SFT (Instruction tuning) & \ref{ch:28}장 & \\
--- & \textbf{평가 (benchmark)} & \textbf{\ref{ch:29}장 \(\leftarrow\) 여기} & OK (학습 아님) \\
4 & Alignment (DPO / GRPO) & \ref{ch:30}·\ref{ch:31}장 & \\
\bottomrule
\end{tabular}
\end{adjustbox}
\end{booktable}

\begin{quote}
능력을 \emph{만든 뒤} (단계 1-3), \emph{얼마나 잘하는지} 를 정량화하는 것이 벤치마크 평가입니다. \emph{정렬 (단계 4)} 로 넘어가기 전, \emph{지금 무엇을 할 수 있는지} 를 먼저 측정합니다.
\end{quote}



\section{변경점 (Diff from \ref{ch:24}--\ref{ch:28}장)}

\begin{booktable}{29장 변경점 요약}{tab:ch29-03}
\begin{adjustbox}{max width=\textwidth}
\begin{tabular}{@{}lll@{}}
\toprule
축 & \ref{ch:24}--\ref{ch:28}장 (학습 장들) & \ref{ch:29}장 (본 장, 평가) \\
\midrule

무엇을 하는가 & \emph{학습} (\inlinecode{Trainer} / \inlinecode{SFTTrainer}) & \textbf{평가만} (추론, \inlinecode{Trainer} 없음) \\
측정 방식 & loss / perplexity & \textbf{task format 별 다른 평가} \\
평가 대상 & 학습 중인 본체 & \textbf{이미 학습된 instruct 모델} \\
출력 형태 & (학습은 출력 없음) & MC 는 \emph{생성 안 함} (log-likelihood), 생성 task 는 토큰 시퀀스 \\
\bottomrule
\end{tabular}
\end{adjustbox}
\end{booktable}

분류 시대 (\ref{ch:01}--\ref{ch:23}장) 의 평가는 단순했습니다 --- 모델이 고정 클래스 중 하나를 고르면 정답 라벨과 비교해 \inlinecode{accuracy} / \inlinecode{F1} 를 냈습니다. \textbf{생성형 LLM 은 자유 토큰 시퀀스를 출력} 하므로, \emph{무엇을 정답으로 보고 어떻게 비교할지} 가 task 마다 다릅니다. 그래서 \emph{하나의 metric} 이 아니라 \emph{여러 평가 방식의 혼합} 이 필요합니다.

\begin{quote}
핵심 전환: \textbf{“정답 라벨과 비교” (분류) \(\to\) “task format 에 맞는 방식으로 측정” (생성)}. 이 장의 §4 표가 그 \emph{3가지 format} 입니다.
\end{quote}



\section{분류 평가 vs 생성 평가 --- 무엇이 달라지나}

\begin{booktable}{29장 분류 평가 vs 생성 평가 - 무엇이 달라지나 표}{tab:ch29-04}
\begin{adjustbox}{max width=\textwidth}
\begin{tabular}{@{}lll@{}}
\toprule
& 분류 (\ref{ch:01}--\ref{ch:23}장) & 생성 LLM (\ref{ch:24}장-) \\
\midrule

metric & \inlinecode{accuracy} / \inlinecode{F1} (정답 라벨 비교) & \textbf{task 마다 다름} \\
출력 & 고정 클래스 (예: 5개 별점) & 자유 토큰 시퀀스 \\
평가 & 한 가지 방식 & \textbf{여러 방식 혼합} \\
예시 & “이 리뷰는 긍정/부정?” \(\to\) 라벨 일치 & “이 수학 문제를 풀어라” \(\to\) 생성 후 답 추출 \\
\bottomrule
\end{tabular}
\end{adjustbox}
\end{booktable}

분류 모델은 \emph{출력 공간이 닫혀 있어} (클래스가 정해져 있어) 비교가 자명합니다. 생성 모델은 \emph{출력 공간이 열려 있어} (어떤 토큰 시퀀스든 나올 수 있어) --- 같은 정답도 표현이 다양하고, 어떤 task 는 \emph{답이 맞았는지} 를 코드로 채점하기조차 어렵습니다.

\begin{center}\rule{0.5\linewidth}{0.5pt}\end{center}

\section{벤치마크 task format 3가지 --- 이 장의 뼈대}

생성형 LLM 벤치마크는 \emph{평가 방식} 에 따라 크게 3가지로 나뉩니다. 이 장은 앞 두 개를 직접 구현하고, 세 번째는 개념으로 다룹니다.

\begin{booktable}{29장 벤치마크 task format 3가지 - 이 장의 뼈대 표}{tab:ch29-05}
\begin{adjustbox}{max width=\textwidth}
\begin{tabular}{@{}lllll@{}}
\toprule
format & 평가 방식 & 생성? & 벤치마크 예 & 본 장 \\
\midrule

\textbf{1 Multiple-choice} & 각 선택지의 \emph{log-likelihood} 를 비교해 argmax & \textbf{생성 안 함} & MMLU, KMMLU, HellaSwag, ARC, KoBEST & \textbf{§2 직접 구현} \\
\textbf{2 Generation + 정답 추출} & 생성한 뒤 정답을 파싱 / 실행해 채점 & 생성함 & GSM8K (수학), HumanEval (코드) & \textbf{§3 시연} \\
\textbf{3 LLM-as-judge} & 다른 (강한) LLM 이 답변의 품질을 채점 & 생성함 & LogicKor, MT-Bench & §6 개념 \\
\bottomrule
\end{tabular}
\end{adjustbox}
\end{booktable}

\begin{quote}
\textbf{핵심 직관}: \emph{객관식 (MC)} 은 모델에게 답을 \emph{쓰게 하지 않습니다}. 대신 \emph{각 선택지를 모델이 얼마나 “그럴듯하다” 고 보는가} (log-likelihood) 를 측정해 가장 높은 것을 고릅니다. 이게 §2 에서 코드로 확인할 \emph{MC 평가의 본질} 입니다. MMLU·KMMLU·HellaSwag 같은 대표 벤치마크 대부분이 이 방식입니다.
\end{quote}



\section{왜 생성형은 정량 평가가 어려운가 --- 쉬운 예제로}

위 표가 \emph{무엇이 다른가} 를 한눈에 보여 줬다면, 여기서는 \emph{왜 어려운가} 를 \textbf{구체적인 예제} 로 손에 잡히게 풀어 봅니다. 핵심 한 문장은 이렇습니다 --- \textbf{분류는 정답이 하나라 채점이 자명하지만, 생성은 정답이 무수히 많아 채점이 주관적입니다.}

\subsection{1 분류 (\ref{ch:01}--\ref{ch:23}장): 정답이 하나, 채점이 자명}

\begin{quote}
입력: \inlinecode{"이\ 영화\ 최고예요"} \(\to\) 출력: \textbf{긍정 / 부정} 둘 중 하나
\end{quote}

정답 라벨이 \emph{정확히 하나} (\inlinecode{긍정}) 입니다. 모델이 \inlinecode{긍정} 이라 하면 맞음, \inlinecode{부정} 이라 하면 틀림. \inlinecode{accuracy} 가 \emph{맞췄나 틀렸나} 로 명확히 정의됩니다. 출력 공간이 \emph{닫혀 있어} (클래스가 정해져 있어) 비교가 기계적입니다.

\subsection{2 생성 (\ref{ch:24}장-): 정답이 무수히 많음, 채점이 주관적}

\begin{quote}
입력: \inlinecode{"건강한\ 식습관\ 3가지\ 알려줘"} \(\to\) 출력: \textbf{무한히 많은 타당한 답}
\end{quote}

\inlinecode{"규칙적인\ 식사,\ 채소\ 섭취,\ 물\ 충분히"} 도 정답이고, \inlinecode{"아침\ 거르지\ 않기,\ 가공식품\ 줄이기,\ 천천히\ 먹기"} 도 정답입니다. \emph{어떤 답이 더 좋은가?} 는 \emph{주관적} 입니다. 정답 라벨 하나와 비교하는 방식이 \emph{통째로 무너집니다}.

\subsection{같은 질문, 여러 타당한 답 --- 능력별로 보면}

\begin{booktable}{29장 같은 질문, 여러 타당한 답 - 능력별로 보면 표}{tab:ch29-06}
\begin{adjustbox}{max width=\textwidth}
\begin{tabular}{@{}lll@{}}
\toprule
능력 & 정답이 하나인가? & 그래서 어떻게 채점? \\
\midrule

\textbf{수학} (GSM8K) & 최종 \emph{숫자} 는 하나지만 \emph{풀이 과정} 은 다양 & 생성 후 \textbf{최종 숫자만 추출} 해 비교 (§3) \\
\textbf{번역·요약} & \emph{여러 표현} 이 다 정답 (\inlinecode{"고양이가\ 잤다"} = \inlinecode{"고양이는\ 잠들었다"}) & BLEU·ROUGE 같은 \emph{n-gram 매칭} --- 그러나 표현이 다르면 정답도 점수가 깎이는 \textbf{한계} \\
\textbf{열린 질문} & \emph{정답 자체가 없음} (식습관 조언, 에세이 ...) & LLM judge 또는 사람 평가 (3) \\
\bottomrule
\end{tabular}
\end{adjustbox}
\end{booktable}

\subsection{부분 정답 문제 --- 이진 채점이 안 됨}

분류는 \emph{맞음 / 틀림} 의 이진입니다. 생성은 그 사이가 넓습니다.

\begin{itemize}
\tightlist
\item
  \emph{절반만 맞은 답}: “건강한 식습관 3가지” 를 물었는데 2가지만 맞게 답함 \(\to\) 0점? 0.67점?
\item
  \emph{형식은 틀렸지만 내용은 맞은 답}: 정답 숫자 \inlinecode{24} 를 \inlinecode{"이십사"} 라고 답함 \(\to\) 단순 문자열 비교로는 \emph{틀림}
\item
  \emph{맞지만 군더더기가 많은 답}: 정답에 불필요한 설명이 섞임 \(\to\) 어디까지 정답?
\end{itemize}

이런 \emph{부분 정답} 은 분류의 라벨 비교로는 표현할 수 없습니다. 그래서 단순 매칭 대신 \emph{task 마다 다른 채점 방식} 이 필요합니다.



\subsection{짧은 코드 --- 왜 exact match 가 부족한가}

같은 질문에 \emph{형식만 다른 두 정답} 을 단순 string match (exact match) 로 채점하면, \textbf{둘 다 “틀림” 으로 나옵니다}. 내용은 맞는데도 말이죠. 아래에서 직접 시연합니다 (모델 없이 문자열만 비교 --- 즉시 실행).



\begin{quote}
위에서 \inlinecode{answer\_a} 는 \emph{숫자 추출} 로는 맞다고 잡히지만, \inlinecode{answer\_b} (\inlinecode{"이십사"}) 는 \emph{숫자 추출조차} 실패합니다 --- 내용은 정답인데도요. \textbf{이것이 생성형 정량 평가의 본질적 어려움} 입니다. 정답을 \emph{인식하는 것 자체} 가 task 마다 다른 규칙을 요구합니다. §2 의 MC 평가가 \emph{생성을 피해} log-likelihood 만 보는 것도, 이 형식 변동 문제를 우회하려는 설계입니다.
\end{quote}



\section{\texorpdfstring{평가는 어쩌다 \emph{하나의 거대한 분야} 가 되었나}{평가는 어쩌다 하나의 거대한 분야 가 되었나}}

여기서 한 걸음 물러나 큰 그림을 봅니다. \textbf{모델 능력이 올라갈수록 평가는 더 어렵고 더 중요해집니다.}

\begin{itemize}
\tightlist
\item
  쉬운 문제는 최신 모델이 \emph{다 풀어 버립니다}. 초창기 벤치마크 (예: 단순 분류·문법 판정) 는 이미 포화 상태입니다.
\item
  그래서 \emph{어려운 능력} (전문 지식, 다단계 추론, 코드, 안전성) 을 구별하려면 \emph{점점 더 정교한 벤치마크} 가 필요합니다. 평가가 모델을 \emph{뒤쫓아} 함께 어려워집니다.
\end{itemize}

그 결과 평가는 \emph{metric 하나} 가 아니라 \textbf{하나의 거대한 분야} 로 자랐습니다.

\begin{itemize}
\tightlist
\item
  능력별 \emph{수십 개} 벤치마크 (지식·추론·수학·코드·진실성·한국어 ...)
\item
  자동 \emph{리더보드} 와 사람 투표 \emph{arena}
\item
  \emph{LLM-as-judge} 로 열린 질문까지 채점
\item
  \emph{안전성·편향·환각} 평가 트랙
\end{itemize}

매주 새 모델과 새 벤치마크가 쏟아지는 지금, \emph{“이 방대한 평가 생태계를 어떻게 항해하는가”} 자체가 실무 역량이 되었습니다.

\begin{quote}
\textbf{상세 평가 부록 안내} --- 이 방대한 생태계를 \emph{실무자 관점} 에서 항해하는 지도를 별도 부록으로 두었습니다: \href{./appendix_eval_landscape.ipynb}{\textbf{\inlinecode{appendix\_eval\_landscape.ipynb}}}. \emph{벤치마크 생태계 지도 · 벤치마크의 진화 · 리더보드 활용 · 평가 도구 비교 · LLM-as-judge · 평가의 함정 · 실무자를 위한 항해 전략} 을 다룹니다. 본 장(§2-§5)가 \emph{평가의 원리} 라면, 부록은 \emph{그 원리를 실무에서 어떻게 쓰는가} 입니다.
\end{quote}



\section{환경 준비}

평가 장이라 학습 라이브러리는 가볍게, 대신 표준 평가 도구 \inlinecode{lm-eval} (lm-evaluation-harness) 를 설치합니다. §2-§4 의 직접 구현은 \inlinecode{transformers} + \inlinecode{datasets} 만으로 동작하고, \inlinecode{lm-eval} 은 §5 (표준 도구 소개) 에서만 씁니다.

\begin{quote}
\inlinecode{lm-eval} 은 의존성이 많아 설치에 1-2분 걸립니다. §5 를 건너뛰려면 \inlinecode{lm-eval} 설치 라인을 빼도 §1-§4 는 그대로 동작합니다.
\end{quote}



\Needspace{11\baselineskip}
\begin{lstlisting}
import re
import random

import numpy as np
import pandas as pd
import torch
from transformers import AutoModelForCausalLM, AutoTokenizer

if torch.cuda.is_available():
    device = torch.device("cuda")
    print(
        f"device : cuda  ("
        f"{torch.cuda.get_device_name(0)})"
    )
elif torch.backends.mps.is_available():
    device = torch.device("mps")
\end{lstlisting}

\noindent\emph{앞 블록은 필요한 라이브러리와 기본 설정을 준비하는 단계입니다. 이어지는 블록에서는 앞에서 만든 중간 값을 다음 계산으로 넘기는 단계로 넘어갑니다.}

\Needspace{11\baselineskip}
\begin{lstlisting}[firstnumber=17]
    print("device : mps  (Apple Silicon)")
else:
    device = torch.device("cpu")
    print(
        "device : cpu  (evaluation will be slow - Colab "
        "T4 recommended)"
    )

USE_FP16 = (device.type == "cuda")
DTYPE = torch.float16 if USE_FP16 else torch.float32

SEED = 0
torch.manual_seed(SEED)
np.random.seed(SEED)
random.seed(SEED)

\end{lstlisting}

\noindent\emph{앞뒤 블록은 모두 앞에서 만든 중간 값을 다음 계산으로 넘기는 단계입니다. 길이를 나누어 같은 흐름을 단계별로 확인합니다.}

\Needspace{5\baselineskip}
\begin{lstlisting}[firstnumber=33]
print(f"torch  : {torch.__version__}")
print(f"fp16   : {USE_FP16}")
\end{lstlisting}



\section{평가 대상 모델 로드}

\textbf{\inlinecode{Qwen/Qwen2.5-0.5B-Instruct}} --- 약 0.5B (494M) 파라미터의 \emph{작은 instruct 모델} 입니다. T4 에서 가볍게 돌고, 한국어·영어를 모두 지원해 한국어 벤치마크에서도 \emph{random 보다 나은} 점수를 냅니다. 작은 모델이라 점수 자체는 낮지만, \emph{평가 파이프라인을 끝까지 돌려보기} 에 적합합니다.

\begin{quote}
비교용으로 \ref{ch:28}장에서 만든 KoGPT2 SFT 모델을 함께 평가할 수도 있습니다. KoGPT2 (125M) 는 너무 작아 대부분의 벤치마크에서 \emph{거의 random} 입니다 --- \textbf{그 자체가 §7 의 교훈}: 작은 모델의 벤치마크 한계. 본 노트북은 Qwen 을 메인으로 진행합니다.
\end{quote}



\Needspace{11\baselineskip}
\begin{lstlisting}
MODEL_NAME = "Qwen/Qwen2.5-0.5B-Instruct"

tokenizer = AutoTokenizer.from_pretrained(MODEL_NAME)
model = AutoModelForCausalLM.from_pretrained(MODEL_NAME, dtype=DTYPE).to(device)
model.eval()

n_params = sum(p.numel() for p in model.parameters())
print(f"model     : {MODEL_NAME}")
print(f"params    : {n_params/1e6:.1f}M")
print(f"vocab     : {tokenizer.vocab_size}")
print(f"eos token : {tokenizer.eos_token!r}")
\end{lstlisting}

\begin{codeRead}
\textbf{1행}의 \inlinecode{MODEL\_NAME = ...}에서는 이후 분석에서 사용할 값을 준비합니다. \textbf{3행}의 \inlinecode{tokenizer = ...}에서는 이후 분석에서 사용할 값을 준비합니다. \textbf{4행}의 \inlinecode{model = ...}에서는 이후 분석에서 사용할 값을 준비합니다.
\end{codeRead}

\noindent\textbf{출력.}
\begin{bookoutputbox}
If the error persists, please let us know by opening an issue on GitHub (ht...
model     : Qwen/Qwen2.5-0.5B-Instruct
params    : 494.0M
vocab     : 151643
eos token : '<|im_end|>'
\end{bookoutputbox}
\par\vspace{0.9em}



\section{Multiple-choice 평가 직접 구현 (핵심)}

객관식 벤치마크 (MMLU, KMMLU, HellaSwag, KoBEST ...) 는 \textbf{모델에게 답을 생성하게 하지 않습니다}. 대신 각 선택지를 \emph{문맥에 이어붙였을 때 모델이 얼마나 그럴듯하게 보는가} --- 즉 \textbf{log-likelihood} 를 계산해 가장 높은 선택지를 정답으로 예측합니다.

\subsection{왜 생성이 아니라 log-likelihood 인가}

\begin{itemize}
\tightlist
\item
  생성하면 \emph{형식 변동} (모델이 “정답은 3번” vs “세 번째” vs 단순히 본문을 이어 씀) 때문에 채점이 불안정합니다.
\item
  log-likelihood 는 \emph{각 선택지에 대한 모델의 확신} 을 직접 수치로 비교 --- 형식에 흔들리지 않고 \emph{재현 가능} 합니다.
\end{itemize}

\subsection{계산 원리}

선택지 \(c\) 의 토큰을 \((t_1, ..., t_k)\) 라 하면, 문맥 (prompt) 뒤에 이어질 확률의 로그는

\begin{equation}
\label{eq:ch29-01}
\log P(c \mid \text{prompt}) = \sum_{i=1}^{k} \log P(t_i \mid \text{prompt}, t_{<i})
\end{equation}

식~\eqref{eq:ch29-01}은 이 절에서 사용하는 핵심 관계를 정리한 것입니다.

모델의 \inlinecode{logits} 에 \inlinecode{log\_softmax} 를 씌워 \emph{각 선택지 토큰 위치의 정답 토큰 log-prob} 를 뽑아 더하면 됩니다. \textbf{\inlinecode{lm-eval-harness} 도 내부에서 이 계산을 합니다} --- 우리는 그 원리를 그대로 구현합니다.



\subsection{log-prob 합 (sum) vs 평균 (mean) --- 길이 정규화}

선택지마다 토큰 수가 다르면 \emph{단순 합} 은 \emph{짧은 선택지} 에 유리합니다 (log-prob 는 음수라, 토큰이 적을수록 합이 덜 깎임). 그래서 \textbf{토큰 수로 나눈 평균 log-prob} (length-normalized) 를 쓰기도 합니다.

\begin{booktable}{29장 log-prob 합 (sum) vs 평균 (mean) - 길이 정규화 표}{tab:ch29-07}
\begin{adjustbox}{max width=\textwidth}
\begin{tabular}{@{}lll@{}}
\toprule
방식 & 식 & 성질 \\
\midrule

sum (\inlinecode{acc}) & \(\sum_i \log P(t_i)\) & 길이가 비슷한 선택지에 적합 \\
mean (\inlinecode{acc\_norm}) & \(\frac{1}{k}\sum_i \log P(t_i)\) & 길이 편향 완화 (HellaSwag 처럼 선택지 길이가 다를 때) \\
\bottomrule
\end{tabular}
\end{adjustbox}
\end{booktable}

\inlinecode{lm-eval-harness} 가 \inlinecode{acc} 와 \inlinecode{acc\_norm} 두 점수를 함께 내는 이유가 이것입니다. 아래에서 둘 다 계산해 비교합니다.



\Needspace{11\baselineskip}
\begin{lstlisting}
@torch.no_grad()
def continuation_logprob(prompt: str, continuation: str):
    '''prompt 뒤에 continuation 이 이어질 때, continuation 토큰들의 log-prob 를
    (sum, mean) 으로 반환. teacher forcing - 생성하지 않고 한 번의 forward 로 계산.'''
    prompt_ids = tokenizer(prompt, return_tensors="pt").input_ids
    full_ids = tokenizer(prompt + continuation, return_tensors="pt").input_ids

    cont_len = max(
        1,
        full_ids.shape[1] - prompt_ids.shape[1],
    )

\end{lstlisting}

\noindent\emph{앞 블록은 텍스트를 모델 입력 텐서로 바꾸는 단계입니다. 이어지는 블록에서는 모델 출력을 확률이나 예측값으로 바꾸는 단계로 넘어갑니다.}

\Needspace{11\baselineskip}
\begin{lstlisting}[firstnumber=13]
    full_ids = full_ids.to(device)
    logits = model(full_ids).logits[0]
    # (T, V)
    log_probs = torch.log_softmax(
        logits.float(),
        dim=-1)  # (T, V,
    )

    target = full_ids[0, 1:]
    # (T-1,)
    pred_lp = log_probs[:-1]                            # (
        T-1,
        V,
    )
    token_lp = pred_lp[torch.arange(
        target.shape[0]),
\end{lstlisting}

\noindent\emph{앞뒤 블록은 모두 모델 출력을 확률이나 예측값으로 바꾸는 단계입니다. 길이를 나누어 같은 흐름을 단계별로 확인합니다.}

\Needspace{11\baselineskip}
\begin{lstlisting}[firstnumber=29]
        target]  # (T-1,,
    )

    cont_lp = token_lp[-cont_len:]
    return cont_lp.sum().item(), cont_lp.mean().item()

def mc_predict(prompt: str, choices: list[str]):
    '''각 선택지의 (sum, mean) log-prob 를 구해 argmax. 두 방식의 예측을 모두 반환.'''
    sums, means = [], []
    for c in choices:
        s, m = continuation_logprob(prompt, c)
        sums.append(s)
        means.append(m)
    return int(np.argmax(sums)), int(np.argmax(means)), sums, means

\end{lstlisting}

\noindent\emph{앞뒤 블록은 모두 모델 출력을 확률이나 예측값으로 바꾸는 단계입니다. 길이를 나누어 같은 흐름을 단계별로 확인합니다.}

\Needspace{11\baselineskip}
\begin{lstlisting}[firstnumber=44]
demo_prompt = "1 더하기 1은 "
demo_choices = ["2입니다.", "3입니다.", "5입니다.", "10입니다."]
pred_sum, pred_mean, sums, means = mc_predict(
    demo_prompt,
    demo_choices,
)
demo_df = pd.DataFrame({
    "choice": demo_choices,
    "logprob_sum": [round(x, 2) for x in sums],
    "logprob_mean": [round(x, 3) for x in means],
})
print(demo_df.to_string(index=False))
print(f"\npredicted (sum)  : {demo_choices[pred_sum]}")
print(f"predicted (mean) : {demo_choices[pred_mean]}")
\end{lstlisting}

\noindent\textbf{출력.}
\begin{bookoutputbox}
choice  logprob_sum  logprob_mean
 2입니다.        -5.01        -1.670
 3입니다.        -5.42        -1.808
 5입니다.        -6.45        -2.151
10입니다.        -7.59        -1.898

predicted (sum)  : 2입니다.
predicted (mean) : 2입니다.
\end{bookoutputbox}
\par\vspace{0.9em}



\subsection{KoBEST HellaSwag --- 4지선다 상식추론}

\textbf{\inlinecode{skt/kobest\_v1}} 의 \inlinecode{hellaswag} task 는 \emph{문맥 (context)} 뒤에 가장 자연스러운 \emph{다음 문장} 을 4개 후보 (\inlinecode{ending\_1..4}) 중에서 고르는 한국어 상식추론 벤치마크입니다. 선택지 길이가 제각각이라 \emph{길이 정규화 (mean)} 효과가 잘 드러납니다. T4 + 시간 제약상 \textbf{test split 의 앞 50문항} 만 평가합니다.



\subsection{KoBEST BoolQ --- 2지선다 (예 / 아니오)}

\inlinecode{boolq} task 는 \emph{본문 (paragraph)} 과 \emph{질문 (question)} 을 주고 \textbf{예 / 아니오} 를 묻습니다. 선택지가 2개라 \emph{random baseline 은 0.5} --- MC 평가를 \emph{가장 단순한 형태} 로 보여줍니다. 여기서는 질문 뒤에 “예” / “아니오” 를 이어붙여 log-prob 를 비교합니다 (label 0 = 아니오, 1 = 예).



\begin{quote}
\textbf{여기까지가 MC 평가의 전부입니다.} 생성은 한 번도 하지 않았습니다 --- 오직 \emph{각 선택지의 log-likelihood} 를 forward 한 번씩으로 계산해 argmax 했을 뿐입니다. 점수가 random 근처라면, 그건 \emph{0.5B 라는 작은 모델의 한계} 입니다 (§7). 같은 코드를 더 큰 모델에 그대로 적용하면 점수가 오릅니다.
\end{quote}



\section{Generation + 정답 추출 평가 (생성 기반)}

두 번째 format 은 \textbf{모델이 실제로 답을 생성} 한 뒤, 그 텍스트에서 \emph{정답을 파싱} 해 채점합니다. GSM8K (초등 수학), HumanEval (코드 실행) 가 대표적입니다. 여기서는 간단한 \textbf{산술 문제 subset} 으로 \emph{생성 \(\to\) 정규식으로 숫자 추출 \(\to\) 정답 비교} 흐름을 시연합니다.

\subsection{few-shot prompt}

모델에게 \emph{예시 몇 개} (few-shot) 를 먼저 보여줘 \emph{답변 형식} 을 유도합니다. 예시 없이 (zero-shot) 던지면 모델이 형식을 못 맞춰 추출이 실패하기 쉽습니다 --- §4 에서 그 차이를 정량으로 봅니다.



\Needspace{11\baselineskip}
\begin{lstlisting}
ARITHMETIC = [
    ("6 곱하기 4는 얼마인가요?", 24),
    ("15 더하기 9는 얼마인가요?", 24),
    ("20 빼기 8은 얼마인가요?", 12),
    ("7 곱하기 7은 얼마인가요?", 49),
    ("100 빼기 37은 얼마인가요?", 63),
    ("13 더하기 28은 얼마인가요?", 41),
]

FEWSHOT = (
    "Q: 사과 3개와 5개를 더하면 몇 개인가요?\nA: 정답은 8입니다.\n\n"
    "Q: 12에서 7을 빼면 얼마인가요?\nA: 정답은 5입니다.\n\n"
)

\end{lstlisting}

\noindent\emph{앞 블록은 앞에서 만든 중간 값을 다음 계산으로 넘기는 단계입니다. 이어지는 블록에서는 텍스트를 모델 입력 텐서로 바꾸는 단계로 넘어갑니다.}

\Needspace{11\baselineskip}
\begin{lstlisting}[firstnumber=15]
def extract_first_int(text: str):
    '''생성 텍스트에서 첫 정수를 추출 (정답 파싱). 없으면 None.'''
    m = re.findall(r"-?\d+", text)
    return int(m[0]) if m else None

@torch.no_grad()
def generate_answer(prompt: str, max_new_tokens: int = 24):
    enc = tokenizer(prompt, return_tensors="pt").to(device)
    out = model.generate(
        **enc,
        max_new_tokens=max_new_tokens,
        do_sample=False,
        pad_token_id=tokenizer.eos_token_id,
    )

\end{lstlisting}

\noindent\emph{앞뒤 블록은 모두 텍스트를 모델 입력 텐서로 바꾸는 단계입니다. 길이를 나누어 같은 흐름을 단계별로 확인합니다.}

\Needspace{11\baselineskip}
\begin{lstlisting}[firstnumber=30]
    return tokenizer.decode(out[0, enc.input_ids.shape[1]:], skip_special_tokens=True)

def eval_generation(problems, shots: str):
    correct = 0
    rows = []
    for q, ans in problems:
        prompt = shots + f"Q: {q}\nA:"
        gen = generate_answer(prompt).strip()
        pred = extract_first_int(gen)
        ok = (pred == ans)
        correct += int(ok)
        rows.append({"question": q, "generated": gen[:30], "pred": pred, "answer": ans, "ok": ok})
    return correct / len(problems), pd.DataFrame(rows)

\end{lstlisting}

\noindent\emph{앞 블록은 텍스트를 모델 입력 텐서로 바꾸는 단계입니다. 이어지는 블록에서는 앞에서 만든 중간 값을 다음 계산으로 넘기는 단계로 넘어갑니다.}

\Needspace{8\baselineskip}
\begin{lstlisting}[firstnumber=44]
acc_gen, gen_df = eval_generation(ARITHMETIC, FEWSHOT)
print(gen_df.to_string(index=False))
print(
    f"\nfew-shot 산술 정확도 : {acc_gen:.3f}  (n="
    f"{len(ARITHMETIC)})"
)
\end{lstlisting}

\begin{codeRead}
\textbf{1--8행}의 \inlinecode{ARITHMETIC = ...}에서는 이후 분석에서 사용할 값을 준비합니다. \textbf{10--13행}의 \inlinecode{FEWSHOT = ...}에서는 이후 분석에서 사용할 값을 준비합니다. \textbf{15행}의 \inlinecode{def extract\_first\_int(tex...}에서는 앞 단계에서 만든 값을 바탕으로 다음 계산을 수행합니다.
\end{codeRead}

\noindent\textbf{출력.}
\begin{bookoutputbox}
question                        generated  pred  answer   ok
6 곱하기 4는 얼마인가요? 정답은 24입니다. \n\n이런 문제들은 어떤 숫자를 곱...
15 더하기 9는 얼마인가요? 정답은 24입니다. \n\n이런 문제들은 어떤 숫자를 더...
20 빼기 8은 얼마인가요? 정답은 12입니다. \n\n이런 문제들은 대부분의 경우, 1...
7 곱하기 7은 얼마인가요? 정답은 49입니다. \n\n이런 문제들은 어떤 숫자를 더...
100 빼기 37은 얼마인가요? 정답은 63입니다. \n\n이런 문제들은 대부분 숫자의...
13 더하기 28은 얼마인가요? 정답은 41입니다. \n\n이런 문제들은 어떤 숫자를...

few-shot 산술 정확도 : 1.000  (n=6)
\end{bookoutputbox}
\par\vspace{0.9em}



\begin{quote}
\textbf{생성 평가의 어려움}: 모델이 정답 숫자를 \emph{맞게 말해도} 형식이 다르면 (\inlinecode{24} vs \inlinecode{이십사} vs 본문에 다른 숫자가 먼저 등장) 정규식 추출이 어긋날 수 있습니다. GSM8K·HumanEval 의 공식 평가가 \emph{정교한 파싱 / 코드 실행} 을 쓰는 이유입니다. MC (§2) 가 \emph{형식에 안 흔들리는} 것과 대조됩니다.
\end{quote}



\section{zero-shot vs few-shot --- in-context learning}

같은 산술 task 를 \textbf{예시 없이 (zero-shot)} 과 \textbf{예시 2개 (2-shot)} 으로 평가해 점수 차이를 봅니다. 모델 가중치는 \emph{전혀 바뀌지 않는데} (학습이 아닙니다), 프롬프트에 예시를 넣는 것만으로 성능이 달라지는 현상이 \textbf{in-context learning} 입니다.

\begin{itemize}
\tightlist
\item
  \textbf{zero-shot}: 예시 없이 문제만. 모델이 \emph{답변 형식} 을 스스로 정해야 함 \(\to\) 추출 실패 잦음
\item
  \textbf{few-shot}: 예시가 \emph{형식 (정답은 N입니다)} 을 보여줘 모델이 그대로 따라 함 \(\to\) 추출 성공률·정확도 상승
\end{itemize}



\begin{quote}
few-shot 이 점수를 올리는 핵심은 \emph{지식을 새로 가르치는 게 아니라} \textbf{답변 형식을 정렬} 시키는 데 있습니다. MMLU 같은 벤치마크가 보통 \emph{5-shot} 으로 보고되는 이유 --- 모델이 \emph{객관식 답 형식} 에 적응하도록.
\end{quote}



\section{\texorpdfstring{\inlinecode{lm-evaluation-harness} 소개 --- 표준 도구}{lm-evaluation-harness 소개 --- 표준 도구}}

§2-§4 에서 직접 구현한 것과 \textbf{정확히 같은 원리} 를, EleutherAI 의 \textbf{\inlinecode{lm-evaluation-harness}} (\inlinecode{lm-eval} 패키지) 가 \emph{표준화된 방식} 으로 수행합니다. 수백 개 벤치마크 (MMLU, HellaSwag, GSM8K, KoBEST, KMMLU ...) 가 \emph{task 정의로 내장} 되어, \emph{프롬프트 포맷 · few-shot 예시 선택 · log-likelihood 계산 · 정규식 추출} 이 모두 통일됩니다. 논문·리더보드의 점수가 이 도구로 측정됩니다.

\subsection{직접 구현 (§2) 과의 관계}

\begin{booktable}{29장 직접 구현 (§2) 과의 관계 표}{tab:ch29-08}
\begin{adjustbox}{max width=\textwidth}
\begin{tabular}{@{}lll@{}}
\toprule
& 직접 구현 (§2) & \inlinecode{lm-eval-harness} \\
\midrule

MC log-likelihood & \inlinecode{continuation\_logprob} 손으로 & 내부에서 동일 계산 (\inlinecode{loglikelihood}) \\
프롬프트 포맷 & 우리가 문자열 조립 & task yaml 에 정의 \\
few-shot 선택 & 우리가 고정 & seed 기반 자동 샘플링 \\
결과 & \inlinecode{acc} 하나 & \inlinecode{acc} + \inlinecode{acc\_norm} + stderr \\
\bottomrule
\end{tabular}
\end{adjustbox}
\end{booktable}

\begin{quote}
\emph{원리는 같고, 표준화·재현성이 다릅니다.} 직접 구현으로 \emph{무슨 일이 일어나는지} 를 이해한 뒤, 실제 보고용 점수는 \inlinecode{lm-eval} 로 내는 것이 일반적인 흐름입니다.
\end{quote}



\subsection{\texorpdfstring{\inlinecode{lm\_eval.simple\_evaluate} API (실행은 선택)}{lm\_eval.simple\_evaluate API (실행은 선택)}}

\inlinecode{lm-eval} 의 파이썬 API 핵심은 \inlinecode{simple\_evaluate} 한 함수입니다. 아래는 \emph{KoBEST BoolQ} 한 task 를 우리 Qwen 모델로 평가하는 코드입니다. \textbf{\inlinecode{lm-eval} 은 버전마다 task 이름·인자가 달라질 수 있어}, 무거우면 실행을 건너뛰고 \emph{사용법} 만 익혀도 됩니다 (직접 구현 §2 가 메인). 셀 상단에서 설치 여부를 확인하고, 설치돼 있을 때만 실행합니다.



\Needspace{11\baselineskip}
\begin{lstlisting}
if HAS_LM_EVAL:
    from lm_eval.models.huggingface import HFLM

    lm = HFLM(
        pretrained=model,
        tokenizer=tokenizer,
        batch_size=8,
    )

    results = lm_eval.simple_evaluate(
        model=lm,
        tasks=["kobest_boolq"],
        num_fewshot=0,
        limit=50,
    )

\end{lstlisting}

\noindent\emph{앞 블록은 필요한 라이브러리와 기본 설정을 준비하는 단계입니다. 이어지는 블록에서는 예측 결과를 지표와 표로 요약하는 단계로 넘어갑니다.}

\Needspace{11\baselineskip}
\begin{lstlisting}[firstnumber=17]
    for task, metrics in results["results"].items():
        print(f"[{task}]")
        for k, v in metrics.items():
            if isinstance(v, float):
                print(f"  {k:14s}: {v:.3f}")
else:
    print(
        "lm-eval 미설치 - §2 의 직접 구현 결과를 표준 "
        "점수로 참고하세요."
    )
\end{lstlisting}



\begin{quote}
위 셀이 버전 문제로 실패하면 (\texttt{task\ \textquotesingle{}kobest\_boolq\textquotesingle{}\ not\ found} 등), \inlinecode{lm\_eval.tasks.TaskManager().all\_tasks} 로 \emph{설치된 버전의 task 이름} 을 확인해 바꾸면 됩니다. \textbf{핵심은 점수 자체가 아니라, §2 의 직접 구현과 \inlinecode{lm-eval} 이 \emph{같은 log-likelihood 원리} 위에 있다는 것} 입니다.
\end{quote}



\section{분야별 벤치마크 지도 --- 무엇을 측정하나}

벤치마크는 \emph{측정하는 능력} 에 따라 분류됩니다. 하나의 점수가 아니라 \emph{여러 능력} 을 \emph{각기 다른 벤치마크} 로 재는 것이 현대 LLM 평가입니다. 아래는 영어·한국어 대표 벤치마크를 능력별로 정리한 지도입니다.

\begin{booktable}{29장 분야별 벤치마크 지도 - 무엇을 측정하나 표}{tab:ch29-09}
\begin{adjustbox}{max width=\textwidth}
\begin{tabular}{@{}llll@{}}
\toprule
측정 능력 & 영어 벤치마크 & 한국어 벤치마크 & format \\
\midrule

\textbf{지식} (전문 분야) & MMLU & KMMLU, HAERAE-Bench & 1 MC \\
\textbf{상식추론} & HellaSwag, ARC & KoBEST (HellaSwag/COPA/BoolQ) & 1 MC \\
\textbf{수학} & GSM8K, MATH & --- (GSM8K 번역본) & 2 생성+추출 \\
\textbf{코드} & HumanEval, MBPP & --- & 2 생성+실행 \\
\textbf{진실성} & TruthfulQA & --- & 1 MC (+ 생성) \\
\textbf{종합 대화/지시} & MT-Bench & LogicKor & 3 LLM-judge \\
\bottomrule
\end{tabular}
\end{adjustbox}
\end{booktable}

\subsection{format 별 다시 보기}

\begin{itemize}
\tightlist
\item
  \textbf{1 MC} (지식·상식·진실성): §2 에서 직접 구현한 \emph{log-likelihood argmax}. 가장 흔하고 재현성 높음
\item
  \textbf{2 생성+추출} (수학·코드): §3 의 \emph{생성 후 파싱/실행}. 채점이 까다롭지만 \emph{실제 풀이 능력} 측정
\item
  \textbf{3 LLM-judge} (대화·지시): 다른 강한 LLM (예: GPT-4) 이 답변을 1-10 점으로 채점. \emph{주관적 품질} 을 측정 --- §7 의 한계 참고
\end{itemize}

\begin{quote}
한 모델을 \emph{제대로 평가} 하려면 이 표 전체를 가로질러야 합니다 --- \emph{지식은 높은데 수학은 약한} 모델이 흔하기 때문입니다. \textbf{한 벤치마크 점수만으로 모델을 판단하면 안 되는} 이유입니다.
\end{quote}



\section{해석 --- 작은 모델의 한계, scaling, 벤치마크 오염}

\subsection{작은 모델의 벤치마크 한계}

§2 에서 Qwen2.5-0.5B 의 KoBEST 점수는 \emph{random 을 살짝 웃도는} 수준이었습니다. \ref{ch:28}장의 KoGPT2 SFT (125M) 를 같은 코드로 평가하면 \emph{거의 정확히 random} 입니다. \textbf{벤치마크는 어느 정도 규모 이상에서만 의미 있는 신호} 를 줍니다 --- 작은 모델은 \emph{지식·추론 용량 자체} 가 부족합니다.

\subsection{scaling 의 필요성}

같은 \inlinecode{continuation\_logprob} 코드를 7B / 70B 모델에 그대로 적용하면 점수가 크게 오릅니다. \emph{코드가 아니라 모델 규모} 가 점수를 만듭니다. 이것이 LLM 에서 \emph{scaling} 이 강조되는 이유 --- 평가 방식은 그대로 두고 모델만 키워도 능력이 따라옵니다.

\subsection{벤치마크 오염 (contamination) 과 한계}

\begin{itemize}
\tightlist
\item
  \textbf{오염}: 벤치마크 문항이 \emph{사전학습 데이터에 섞여} 들어가면, 모델이 \emph{외워서} 맞히는 것일 수 있습니다 (실제 능력 과대평가). 새 모델일수록 의심해야 합니다.
\item
  \textbf{단일 벤치마크의 위험}: MMLU 만 높은 모델이 실제 대화는 형편없을 수 있습니다. \emph{여러 능력 (§6)} 을 가로질러 봐야 합니다.
\item
  \textbf{format 편향}: MC 점수가 높다고 \emph{실제 생성 품질} 이 좋은 건 아닙니다 --- MC 는 \emph{고르기만} 하면 되지 \emph{쓰지} 는 않으니까요.
\end{itemize}

\begin{quote}
평가는 \emph{모델을 만드는 것만큼 어렵습니다}. 좋은 평가 없이는 \emph{무엇이 나아졌는지} 알 수 없고, 다음 단계 (alignment) 를 향한 방향도 잡을 수 없습니다.
\end{quote}



\section{체크포인트 질문}

\begin{enumerate}
\def\labelenumi{\arabic{enumi}.}
\tightlist
\item
  분류 평가 (\ref{ch:01}--\ref{ch:23}장) 와 생성 LLM 평가가 \emph{근본적으로 다른} 이유는 무엇인가요? metric 이 하나가 아닌 까닭을 출력 공간 관점에서 설명해 보세요.
\item
  Multiple-choice 벤치마크는 왜 \emph{생성하지 않고} log-likelihood 만 계산하나요? 생성 방식 대비 어떤 장점이 있나요?
\item
  log-prob 의 \emph{합 (sum)} 과 \emph{평균 (mean)} 중, 선택지 길이가 제각각인 HellaSwag 에서 어느 쪽이 길이 편향에 강한가요? 그 이유는?
\item
  few-shot 예시가 모델 가중치를 바꾸지 않는데도 점수를 올리는 현상의 이름은 무엇이며, 무엇을 “정렬” 시키기 때문인가요?
\item
  한 모델을 \emph{하나의 벤치마크 점수} 로만 판단하면 안 되는 이유를 §6·§7 을 근거로 두 가지 이상 들어 보세요.
\end{enumerate}



\section{FAQ}

\begin{faqBox}{Q1. MC 벤치마크는 왜 답을 생성하게 하지 않고 log-likelihood 를 보나요?}
생성하면 \emph{형식 변동} 때문에 채점이 불안정합니다 --- 모델이 “정답은 2번”, “두 번째”, 혹은 단순히 본문을 이어 쓸 수도 있습니다. log-likelihood 는 \emph{각 선택지에 대한 모델의 확신} 을 직접 수치로 비교하므로 형식에 흔들리지 않고 \emph{재현 가능} 합니다. 또한 forward 한 번이면 되어 \emph{생성 (autoregressive decode) 보다 빠릅니다}. §2 의 \inlinecode{continuation\_logprob} 이 그 계산의 전부입니다.

\end{faqBox}

\begin{faqBox}{Q2. log-prob 를 토큰 수로 나누는 (mean) 정규화는 왜 필요한가요?}
log-prob 는 음수라, \emph{토큰이 많은 선택지일수록 합이 더 깎입니다}. 그러면 \emph{짧은 선택지} 가 부당하게 유리해집니다. 토큰 수로 나눠 \emph{토큰당 평균 log-prob} 를 보면 이 길이 편향이 완화됩니다. HellaSwag 처럼 선택지 길이가 크게 다른 task 에서 \inlinecode{acc\_norm} (mean) 이 \inlinecode{acc} (sum) 보다 흔히 더 높게 나오는 이유입니다.

\Needspace{5\baselineskip}
\begin{lstlisting}
score_sum  = cont_lp.sum()
score_mean = cont_lp.sum() / cont_len
\end{lstlisting}

\end{faqBox}

\begin{faqBox}{Q4. LLM-as-judge (3) 의 장단점은 무엇인가요?}

\begin{itemize}
\tightlist
\item
  \textbf{장점}: \emph{주관적 품질} (유창함, 도움됨, 안전성) 을 측정할 수 있습니다. MC·정규식으로는 잴 수 없는 \emph{자유 생성} 의 품질을 사람 대신 강한 LLM 이 채점합니다 (MT-Bench, LogicKor).
\item
  \textbf{단점}: judge 모델의 \emph{편향} 을 물려받습니다 (긴 답변·자기 출력 선호, 위치 편향). 비용도 듭니다 (judge 호출). 그래서 \emph{MC·생성 평가와 병행} 하고, judge 결과는 절대 점수보다 \emph{상대 비교} 로 보는 것이 안전합니다.
\end{itemize}

\end{faqBox}

\compactFAQMore{4}

\begin{previewBox}{미리보기: 다음 장}

\textbf{\ref{ch:30}장. DPO --- 사람 선호로 정렬 (Alignment, 학습 단계 4)}

\begin{itemize}
\tightlist
\item
  본 장 (벤치마크) 가 \emph{능력 측정} 이었다면, \ref{ch:30}--\ref{ch:31}장의 \textbf{alignment} 는 \emph{선호·안전성 정렬} 입니다. \emph{얼마나 잘하는가} 에서 \emph{얼마나 사람이 원하는 방식으로 하는가} 로.
\item
  \textbf{DPO (Direct Preference Optimization)}: \ref{ch:28}장의 SFT 모델을 \emph{preference 데이터 (chosen / rejected 쌍)} 로 정렬. 사람이 \emph{더 선호하는} 답변의 확률을 올리고 덜 선호하는 답변의 확률을 내립니다.
\item
  \textbf{\inlinecode{labels\ =\ -100} thread 가 DPO 에서도 이어집니다} --- chosen / rejected \emph{둘 다 response 부분만} log-prob 를 계산해 비교합니다. 본 장 §2 의 \inlinecode{continuation\_logprob} 과 \emph{정확히 같은 log-prob 계산} 이 DPO loss 의 핵심 재료입니다.
\end{itemize}

\textbf{Phase 4 GPT 시대 흐름 --- 본 장의 위치}:

\begin{booktable}{29장 FAQ 표}{tab:ch29-10}
\begin{adjustbox}{max width=\textwidth}
\begin{tabular}{@{}lll@{}}
\toprule
장 & 단계 & 하는 일 \\
\midrule

\ref{ch:24}·\ref{ch:26}장 & 1 (pretraining) & 본체를 만든다 \\
\ref{ch:25}·\ref{ch:27}장 & 2 (continual pretraining) & 도메인에 적응 \\
\ref{ch:28}장 & 3 (SFT) & 지시를 따르게 \\
\textbf{\ref{ch:29}장 \(\leftarrow\) 여기} & \textbf{--- (평가)} & \textbf{무엇을 할 수 있는지 측정} \\
\ref{ch:30}·\ref{ch:31}장 & 4 (alignment) & 사람 선호로 정렬 \\
\bottomrule
\end{tabular}
\end{adjustbox}
\end{booktable}

\begin{quote}
\textbf{변하는 축} (\ref{ch:28}장 \(\to\) \ref{ch:29}장): \emph{학습 \(\to\) 평가}. 모델·loss 가 아니라 \emph{측정 방식} 이 주제입니다. §2 의 MC log-likelihood 가 \ref{ch:30}장 DPO 의 chosen/rejected log-prob 비교로 \emph{재료 그대로} 이어집니다.
\end{quote}
\end{previewBox}



\Needspace{18\baselineskip}

\section{온라인 부록: 생성형 LLM 평가 항해 가이드}

\onlineAppendixLink{생성형 LLM 평가 항해 가이드}{벤치마크 생태계, 리더보드, 평가 도구, LLM-as-judge와 사람 평가의 편향을 한 흐름으로 정리합니다. 공개 점수는 참고 지표로 쓰고 실제 배포 판단은 use-case 맞춤 평가셋으로 내린다는 원칙이 부록의 결론입니다.}{https://colab.research.google.com/github/yoon-gu/neuqes-101/blob/master/29_benchmark_eval/appendix_eval_landscape.ipynb}